% =====================================================================
%  5.7 INTRODUCCION
%  Rubrica (Doc1): para "Excelente" se exige informacion de calidad,
%  que SE PRESENTE EL PROBLEMA y que se justifique TOTALMENTE el trabajo
%  realizado, indicando como se tratara el problema.
%
%  OBLIGATORIO si el TFG integra trabajos de otras materias: hay que
%  especificar aqui esos trabajos y esas materias.
% =====================================================================

\chapter{Introducción}
\label{cap:introduccion}

\section{Contexto del problema}

En las personas mayores de 65 años, la interacción entre medicamentos y su relación con
enfermedades aumenta el riesgo de reacciones adversas~\cite{maher2014polypharmacy}.
Ante esta situación se emplean criterios de cribado como STOPP/START, descritos
en la sección siguiente~\cite{omahony2023stopp}.

\section{Los criterios STOPP/START}

Los criterios STOPP/START (\textit{Screening Tool of Older Persons' Prescriptions} /
\textit{Screening Tool to Alert to Right Treatment}) constituyen una herramienta de
cribado clínico orientada a la revisión de la medicación en personas mayores de 65
años~\cite{omahony2023stopp}. Su objetivo es detectar, por un lado, prescripciones
potencialmente inapropiadas (STOPP) y, por otro, omisiones de tratamientos indicados
(START).

La versión vigente es la tercera (STOPP/START v3), publicada en 2023, que recoge
133 criterios STOPP y 57 criterios START, hasta un total de 190~\cite{omahony2023stopp}.
En la práctica clínica habitual, el profesional debe consultar este listado
---habitualmente distribuido en formato PDF~\cite{sacyl2023stopp}--- e interpretarlo
manualmente sobre el caso del paciente.

\section{Justificación del trabajo}
\label{sec:justificacion}

La aplicación sistemática de los 190 criterios STOPP/START v3 resulta laboriosa si se
realiza de forma manual: el profesional debe contrastar el listado completo con la
medicación y los diagnósticos de cada paciente, algo poco viable en consultas con
elevada carga asistencial. Este TFG responde a esa necesidad con una herramienta de
apoyo a la decisión que automatiza ese contraste y reduce el esfuerzo de la revisión,
sin sustituir el juicio clínico. El enfoque concreto se describe a continuación.

\section{Cómo se aborda el problema}
\label{sec:abordaje}

Para abordar el problema descrito en la~\cref{sec:justificacion}, se
desarrolla una aplicación web que evalúa de forma automática los criterios
sobre el caso introducido por el profesional. El núcleo es un motor de reglas:
una regla por cada criterio, aplicada a la medicación y a los diagnósticos,
de modo que los avisos aplicables aparecen de forma inmediata.
Además, se deshabilitan opciones que no pueden activar ningún criterio con la
selección actual y se destacan las relaciones entre medicamentos y
diagnósticos. El detalle técnico se desarrolla en los capítulos de
arquitectura y diseño.

\section{Integración de trabajos previos}
\label{sec:trabajos-previos}

Este TFG no integra trabajos ni materiales procedentes de prácticas u otras
asignaturas del grado. Sí se aplican, de forma indirecta, conceptos adquiridos
a lo largo de la carrera ---por ejemplo, ingeniería del software o diseño de
interfaces---, pero el sistema se ha
concebido y desarrollado de manera específica para este trabajo.